\chapter{Preface}
\markboth{Preface}{Preface}

\section*{Purpose of This Manual}

This manual tells you how to build, set up, and operate a pen plotter that
uses ClearHPGL. It tells you how to install the firmware on a Teknic
ClearCore, how to connect the servos, the pen lift, and the emergency stop,
and how to set the board for your machine. It also tells you how to plot from
the ClearHPGL desktop app, from AutoCAD, and from other software.

ClearHPGL is firmware for the ClearCore. With it, the ClearCore and two
ClearPath servos operate as an HP-GL pen plotter. A CAD program sends a
drawing through a standard plotter driver, and the machine draws it. The
optional desktop app shows a preview of the drawing, sends jobs, and shows
the pen as it moves.

\section*{Audience}

This manual is for persons who build and operate the machine. You must know
how to connect motors and switches to an industrial controller, and how to
use a serial terminal and a Windows computer. Knowledge of HP-GL helps, but
it is not necessary. The glossary defines the terms that this manual uses.

\section*{How to Use This Manual}

Table~\ref{tab:roadmap} shows where to start for each task.

\begin{table}[htbp]
\centering
\caption{Where to start}\label{tab:roadmap}
\begin{tabularx}{\linewidth}{@{}P{0.46\linewidth}L@{}}
\toprule
\textbf{Task} & \textbf{Read} \\
\midrule
Learn what ClearHPGL does and does not do & Chapter~\ref{ch:intro} \\
Operate the machine safely & Chapter~\ref{ch:safety} \\
Connect the servos, the pen lift, and the emergency stop & Chapter~\ref{ch:hardware} \\
Install or update the firmware & Chapter~\ref{ch:firmware} \\
Set the board for your machine and do the first moves & Chapter~\ref{ch:setup} \\
Plot from the desktop app & Chapter~\ref{ch:app} \\
Plot from AutoCAD or other software & Chapter~\ref{ch:cad} \\
Understand how the board processes a job & Chapter~\ref{ch:jobs} \\
Find the cause of a fault & Chapter~\ref{ch:trouble} \\
Look up a command, a setting, or an HP-GL instruction & Appendixes~\ref{app:console}--\ref{app:network} \\
\bottomrule
\end{tabularx}
\end{table}

\section*{Conventions}

\subsection*{Language}

The text of this manual follows ASD-STE100, \textit{Simplified Technical
English}.\footnote{ASD, \textit{ASD-STE100: Simplified Technical English},
issue~8 (Brussels: ASD, 2021).} STE makes the text easy to read and easy to
translate. These are the rules that you will see most:

\begin{itemize}
  \item A procedure step contains one instruction. Each step starts with a
        verb in the imperative, for example ``Set'' or ``Do not.''
  \item A sentence in a procedure has no more than 20 words. A sentence in
        a description has no more than 25 words.
  \item The text uses the active voice and the simple tenses.
  \item A word has one meaning. For example, ``enable'' always means ``supply
        power to the motor windings,'' and ``plot'' always means ``draw a job
        on the machine.''
  \item A warning or a caution comes before the step that it applies to.
\end{itemize}

Technical names, such as setting names, console commands, HP-GL
instructions, and file names, are not changed.

\subsection*{Style}

The apparatus of this manual follows \textit{The Chicago Manual of
Style}.\footnote{\textit{The Chicago Manual of Style}, 18th ed. (Chicago:
University of Chicago Press, 2024).} Chapter and section titles use headline
style. Lists of three or more items use the serial comma. Ranges use an en
dash, for example IO-0--IO-5. The manual uses the alternative rule for numbers
in Chicago 9.3: words for zero through nine in running text, and numerals for
10 and above. Numerals are used for all quantities with a unit, for setting
values, and for code. Notes are footnotes in the Chicago notes-bibliography
style, and the bibliography is at the end of the manual.

\subsection*{Safety Notices}

\begin{warning}
A warning tells you about a condition that can cause injury or death. It
also tells you about a condition that can make machinery move when you do
not expect it.
\end{warning}

\begin{caution}
A caution tells you about a condition that can cause damage to equipment,
to the material on the bed, or to the settings of the board.
\end{caution}

\begin{note}
\notelabel A note gives information that helps you do a task. It does not
tell you about a hazard.
\end{note}

\subsection*{Typography}

\begin{table}[htbp]
\centering
\caption{Typographic conventions}\label{tab:typo}
\begin{tabularx}{\linewidth}{@{}P{0.32\linewidth}L@{}}
\toprule
\textbf{Example} & \textbf{Meaning} \\
\midrule
\cmd{set width_mm 600} & A console command, an HP-GL instruction, or text that you type. \\
\key{draw_mm_s} & A board setting. Appendix~\ref{app:settings} lists all settings. \\
\file{Firmware/build.ps1} & A file or a directory, relative to the root of the repository. \\
\menu{Plotter settings} & A menu, a button, a field, or a check box in the desktop app. \\
\conn{M0}, \conn{DI-6} & A connector on the ClearCore. \\
\term{hard clip} & A term that the glossary defines, at its first use. \\
\bottomrule
\end{tabularx}
\end{table}

All lengths are in millimeters (mm) and all speeds are in millimeters per
second (mm/s), unless the text gives a different unit. HP-GL uses
\term{plotter units}: one plotter unit is \SI{0.025}{mm}, so there are
40 plotter units in \SI{1}{mm}.

\section*{Related Documents}

The file \file{Firmware/PROTOCOL.md} in the repository is the technical
reference for the firmware. It lists the links, the HP-GL instructions, the
console commands, and the settings. This manual contains that information.
If this manual and \file{PROTOCOL.md} are different, \file{PROTOCOL.md}
agrees with the firmware. Use it, and report the difference.

Teknic supplies the manuals for the ClearCore and for the ClearPath servos.
Use them for the electrical data of the hardware: power supplies,
connector pinouts, and the setup of each servo. The bibliography lists them.
